\documentclass[11pt]{article}
\usepackage[margin=1in]{geometry}
\usepackage{amsmath,amssymb,amsthm,mathtools}
\usepackage{graphicx}
\usepackage{microtype}
\usepackage[hidelinks]{hyperref}

\newtheorem{theorem}{Theorem}
\newtheorem{lemma}{Lemma}
\newcommand{\Q}{\mathbb{Q}}
\newcommand{\C}{\mathbb{C}}
\newcommand{\trdeg}{\operatorname{trdeg}}

\title{Algebraic Independence Forces Simultaneous Nonvanishing\\
\large A full answer to the second final-comments conjecture of arXiv:1502.01049}
\author{Candidate solution packet --- run \texttt{fa\_banach\_001}, lane 0}
\date{27 August 2026}

\begin{document}
\maketitle

\begin{center}
\fbox{\parbox{0.91\textwidth}{\textbf{Status.} Candidate full affirmative
answer, likely valid, subject to human review. The result answers only the
second conjecture in the source paper's final paragraph.}}
\end{center}

\section{Source and exact question}

Philippe Ryckelynck and Laurent Smoch study matricial polynomial functions and
a convolution-type functional equation in \cite{RS}. In their final comments
(source PDF page 15), they conjecture a transcendence-degree variant of their
Theorem 4.1. The source passage is reproduced below.

\begin{figure}[ht]
  \centering
  \includegraphics[width=\textwidth]{figures/open_problem_crop.png}
  \caption{Full-width crop of source PDF page 15. The target is the second
  conjecture, beginning ``Second, we conjecture.''}
\end{figure}

The conjecture is read with the hypotheses inherited from the phrase
``variant of Theorem 4.1'': one is given finitely many \emph{generic}
matricial polynomial functions
\[
 \psi^{(k)}(\zeta,\xi)
 =\sum_m \kappa^{(k)}_m
   \zeta^{m_1}\xi^{m_2}\zeta^{m_3}\xi^{m_4}\cdots,
 \qquad 1\leq k\leq \ell.
\]
If all coefficients are rational, or more generally algebraic over \(\Q\),
and the \(2d^2\) entries of \(\zeta,\xi\in M_d(\C)\) generate a subfield of
\(\C\) of transcendence degree \(2d^2\), must every
\(\psi^{(k)}_{ij}(\zeta,\xi)\) and every
\(\det\psi^{(k)}(\zeta,\xi)\) be nonzero?

The inherited genericity is essential: with no nonzero/generic hypothesis,
the zero polynomial is an immediate counterexample. The source's definition
of genericity is exactly the one used in Theorem 4.1 and its proof
\cite[pp.~6--7]{RS}.

\section{Definitions and proof intuition}

After expanding matrix products, each scalar entry of
\(\psi^{(k)}(\zeta,\xi)\) is an ordinary commutative polynomial in the
\(2d^2\) scalar entries of \((\zeta,\xi)\); so is its determinant. The proof
of source Theorem 4.1 shows that genericity makes every one of these formal
polynomials nonzero.

The remaining step is field-theoretic. A tuple of \(n\) elements generating
a field of transcendence degree \(n\) is algebraically independent. Moreover,
algebraically independent elements remain so after an algebraic extension of
the base field. Here \(n=2d^2\), and the finitely many algebraic coefficients
generate just such an algebraic extension of \(\Q\). Thus none of the source
paper's nonzero polynomials can vanish at the prescribed matrix pair.

\section{Formal result}

\begin{lemma}[Algebraic base change]
Let \(K/F\) be an algebraic field extension and let
\(x=(x_1,\dots,x_n)\) lie in a common overfield. If
\(\trdeg_F F(x)=n\), then \(x_1,\dots,x_n\) are algebraically independent
over \(K\).
\end{lemma}

\begin{proof}
Since \(K/F\) is algebraic, adjoining \(K\) does not change transcendence
degree:
\[
 \trdeg_K K(x)=\trdeg_F F(x)=n.
\]
If the \(n\) elements \(x_1,\dots,x_n\) were algebraically dependent over
\(K\), a transcendence basis could be chosen from fewer than \(n\) of them,
forcing \(\trdeg_K K(x)<n\), a contradiction.
\end{proof}

\begin{theorem}[Affirmative answer to the second conjecture]
Let \(\psi^{(1)},\dots,\psi^{(\ell)}\) be generic matricial polynomial
functions in the sense of Ryckelynck--Smoch Theorem 4.1. Suppose every
coefficient \(\kappa^{(k)}_m\) is algebraic over \(\Q\). Let
\(\zeta,\xi\in M_d(\C)\), and suppose the field generated over \(\Q\) by
their \(2d^2\) scalar entries has transcendence degree \(2d^2\). Then, for
all \(1\leq k\leq\ell\) and \(1\leq i,j\leq d\),
\[
 \psi^{(k)}_{ij}(\zeta,\xi)\ne0,
 \qquad
 \det\!\bigl(\psi^{(k)}(\zeta,\xi)\bigr)\ne0.
\]
In particular, the conjecture holds for rational coefficients.
\end{theorem}

\begin{proof}
List the entries of \((\zeta,\xi)\) as
\(x=(x_1,\dots,x_n)\), where \(n=2d^2\), and set \(E=\Q(x)\).
The hypothesis \(\trdeg_\Q E=n\) implies that the \(n\)-tuple \(x\) is
algebraically independent over \(\Q\): any algebraic dependence would reduce
the maximal size of an algebraically independent subtuple below \(n\).

Let \(K\) be the field generated over \(\Q\) by all coefficients occurring in
the finite family \(\psi^{(1)},\dots,\psi^{(\ell)}\). There are only finitely
many such coefficients, and each is algebraic over \(\Q\), so \(K/\Q\) is a
finite algebraic extension. By the lemma, \(x_1,\dots,x_n\) are algebraically
independent over \(K\).

For every \(i,j,k\), expansion of the matrix words gives an ordinary
polynomial
\[
 P^{(k)}_{ij}(X_1,\dots,X_n)\in K[X_1,\dots,X_n]
\]
whose value at \(x\) is \(\psi^{(k)}_{ij}(\zeta,\xi)\). Likewise,
\[
 D^{(k)}(X_1,\dots,X_n)
 :=\det\bigl(P^{(k)}_{ij}(X_1,\dots,X_n)\bigr)_{i,j=1}^d
 \in K[X_1,\dots,X_n].
\]
The proof of source Theorem 4.1 establishes, from the source's genericity
conditions, that every \(P^{(k)}_{ij}\) and every \(D^{(k)}\) is a nonzero
formal polynomial \cite[pp.~6--7]{RS}. (For determinants, the source
specializes \(\zeta=\rho I_d\), \(\xi=\sigma I_d\), obtaining the \(d\)-th
power of a nonzero scalar polynomial; its upper-triangular specialization and
conjugation argument handles all entries.)

Algebraic independence of \(x\) over \(K\) means precisely that the evaluation
homomorphism
\[
 K[X_1,\dots,X_n]\longrightarrow K(x_1,\dots,x_n),
 \qquad P\longmapsto P(x),
\]
is injective. Therefore no \(P^{(k)}_{ij}\) and no \(D^{(k)}\) can evaluate
to zero. These evaluations are exactly the quantities asserted to be
nonzero, completing the proof.
\end{proof}

\section{Verification, scope, and novelty}

The proof uses no analytic regularity, recurrence theory, or computational
assumption. It requires only (i) the nonzero-formal-polynomial conclusion
proved inside source Theorem 4.1 and (ii) invariance of transcendence degree
under algebraic base extension. The conclusion is simultaneous for every
member of the source's finite family.

The result does \emph{not} address the source's first pseudo-periodicity
problem or its third conjecture on open-dense genericity for arbitrary
\(N\geq2\). It also does not support a reading that drops genericity; that
reading is false for \(\psi=0\).

A bounded search on 27 August 2026 covered the run indexes, the exact arXiv id,
title and authors, the exact transcendence-degree phrase, and close keyword
variants. It found the source and journal/HAL copies but no separate paper
answering this conjecture. Novelty confidence is moderate rather than
exhaustive; no full MathSciNet/zbMATH citation-graph search was available.

\paragraph{Human-review focus.}
Verify that the phrase ``variant of Theorem 4.1'' retains the genericity
hypothesis, and check the algebraic base-change lemma. Under that natural and
necessary reading, the argument is complete.

\begin{thebibliography}{9}
\bibitem{RS}
P. Ryckelynck and L. Smoch,
\emph{On the functional equation \(\alpha\mathbf{u}+\mathcal{C}\star
(\chi\mathbf{u})=\mathbf{f}\)},
North-Western European Journal of Mathematics \textbf{4} (2018), 1--19;
arXiv:1502.01049.
\end{thebibliography}

\end{document}
